\section*{Resumen}
\justifying
RESUMEN PARA DIFUSIÓN: en catálogo de Biblioteca; en eventos y publicaciones de
la Secretaría de Extensión Universitaria y en catálogo colectivo de tesis del Acuerdo de Bibliotecas Universitarias de Córdoba (ABUC).
QUE ES Y COMO SE DEBE HACER un RESUMEN Y las PALABRAS CLAVES
“El abstract debería ser considerado como una miniversión del artículo” (Day 1991).Describe los objetivos del estudio, la metodología usada, los resultados principales del trabajo y sus conclusiones fundamentales. Un abstract tiene típicamente uno o dos párrafos y menos de 200 palabras y debe permitir a los lectores identificar el  contenido básico deI documento rápida y fielmente, con el fin de determinar la relevancia del mismo para sus intereses y, por lo tanto, para decidir si necesitan leer el documento en su Totalidad. 

\noindent \textbf{Palabras clave:}
Las palabras Claves propuestas por el Autor, son para la indexación del documento.

Comentario: El resumen debe constar de 2 párrafos. El primero para contextualizar indicando el ¿Qué problema se resolverá?, ¿Porqué?¿A quienes beneficia? y ¿Dónde?. El segundo descriptivo para indicar ¿Qué se ha hecho?¿Para qué?¿Cómo y Con qué? ¿Qué resultados se obtuvo y Qué se recomienda? La suma de estos párrafos no deberá exceder las 200 palabras.
Para contar las palabras usar la opción de Word: Herramientas/Contar palabras.